\section{Conclusion}
\label{sec:conclusion}

% FINAL-RUN DRAFTING NOTES (reviewed 2026-09-21)
% Write findings and policy implications after Sections 4 and 5. Budget: 800 words
% total including the interpretation and future-work prose already present.
% Answer each RQ briefly, stating failures and uncertainty as well as agreement.
% Distinguish fitted baseline agreement, unfitted checks and counterfactual predictions.
% State the outcome of the nonlinearity test as observed, including a negative result
% if supported. Do not presume either a smooth response or a null visibility effect.
% Contributions: the two-lender-class model, reusable survey-derived population, and
% affordability-channel policy comparison. Bound novelty to the literature reviewed.
% Ground policy interpretation in retained stacking and scenario evidence; null
% effects alone do not establish why an intervention has little effect.
% Keep the scoring channel outside the claimed result and avoid numerical thresholds.
% After results are written, compress the introduction's preview and write the abstract.

\subsection{Summary of Findings}
\label{sec:summary}

\subsection{Contributions}
\label{sec:contributions}

\subsection{Policy Implications}
\label{sec:policy}

\label{sec:lim-reading}
The experiments estimate changes within a model of households under 2017 conditions.
Their policy interpretation depends on the assumptions and sensitivity checks reported alongside
each result. Holding parameters fixed between scenarios isolates the modelled effect, but does
not remove uncertainty about those parameters. In particular, the bureau-visibility comparison
addresses only the affordability channel of the reporting requirement.

\subsection{Future Work}
\label{sec:future-work}

Four extensions address the model's main limitations.

\textbf{Credit scoring.} A rule linking reported \ac{BNPL} behaviour to credit scores, loan
pricing and refusals would allow the model to examine the reporting channel it currently omits
(Section~\ref{sec:asymmetry}). This would require evidence on how bureaus and lenders use the data.

\textbf{Feedback and contagion.} A lender that restricts credit after losses, or a link from
aggregate default to employment, would allow distress to spread through channels currently
excluded (Table~\ref{tab:design-concepts}). Peer effects on consumption, as in Cardaci
\cite{cardaci2018inequality}, would also extend the present adoption mechanism.

\textbf{Adoption parameters.} The \ac{CFPB} stacking shares could help constrain
$q_{\text{base}}$, subject to differences between the United States and South African markets.
If used for calibration, those shares would cease to be a separate plausibility check.

\textbf{South African behavioural evidence.} Local estimates of minimum-payment anchoring,
present bias and peer influence would reduce reliance on imported parameters.
